% Компилятор: pdfLaTeX, XeLaTeX или LuaLaTeX. Для ссылок нужны два прохода.
% Для Times New Roman выберите XeLaTeX/LuaLaTeX и установите этот шрифт.
\documentclass[a4paper,14pt]{extarticle}
\usepackage{iftex}
\ifPDFTeX
  \usepackage[utf8]{inputenc}
  \usepackage[T2A]{fontenc}
  % Times-совместимый шрифт с кириллицей для pdfLaTeX.
  \usepackage{tempora}
  \PackageWarningNoLine{lab-template}{pdfLaTeX uses Tempora, not Times New Roman.
    For Times New Roman select XeLaTeX or LuaLaTeX}
\else
  \usepackage{fontspec}
  % Подключаем файлы напрямую: это исключает сбои поиска системных шрифтов.
  % Файлы Times New Roman можно положить рядом с main.tex.
  \IfFileExists{times.ttf}{%
    \setmainfont{times.ttf}[BoldFont=timesbd.ttf,
      ItalicFont=timesi.ttf,BoldItalicFont=timesbi.ttf]%
  }{%
    \setmainfont{Tempora-Regular.otf}[BoldFont=Tempora-Bold.otf,
      ItalicFont=Tempora-Italic.otf,BoldItalicFont=Tempora-BoldItalic.otf]%
    \PackageWarningNoLine{lab-template}{Times New Roman files are missing.
      Tempora is used. Add times.ttf, timesbd.ttf, timesi.ttf and timesbi.ttf
      for Times New Roman}%
  }
\fi
% У extarticle опция 14pt задаёт 14.4pt; уточняем основной кегль до 14pt.
\makeatletter
\renewcommand\normalsize{\@setfontsize\normalsize{14pt}{17pt}}
\makeatother
\normalsize
\usepackage[russian]{babel}

% false — поля и абзац по ГОСТ; true — параметры из приложенного DOCX.
\newif\ifDepartmentStyle
\DepartmentStylefalse
% Для короткого отчёта можно отключить новую страницу перед разделом.
\newif\ifNewSectionPage
\NewSectionPagetrue

% Шрифт выбирается выше в зависимости от компилятора и доступных гарнитур.
\usepackage{geometry}
\ifDepartmentStyle
  \geometry{left=30mm,right=10mm,top=20mm,bottom=20mm,footskip=10mm}
  \newcommand{\ParagraphIndent}{10mm}
\else
  \geometry{left=30mm,right=15mm,top=20mm,bottom=20mm,footskip=10mm}
  \newcommand{\ParagraphIndent}{12.5mm}
\fi
\usepackage{setspace}
\onehalfspacing
\setlength{\parindent}{\ParagraphIndent}
\setlength{\parskip}{0pt}
\usepackage{indentfirst}
\usepackage{amsmath}
\usepackage{graphicx}
\usepackage{array,tabularx}
\usepackage{caption}
\DeclareCaptionLabelSeparator{gostdash}{\space—\space}
\DeclareCaptionLabelSeparator{departmentdot}{.\space}
\captionsetup{font=normalsize,labelsep=gostdash}
\captionsetup[figure]{name=Рисунок,position=bottom,justification=centering}
\captionsetup[table]{name=Таблица,position=top,justification=raggedright,singlelinecheck=false}
\ifDepartmentStyle
  \captionsetup[table]{labelsep=departmentdot}
\fi
\usepackage{titlesec}
\titleformat{\section}{\normalfont\bfseries\hyphenpenalty=10000}{\thesection}{1em}{}
\titleformat{\subsection}{\normalfont\bfseries\hyphenpenalty=10000}{\thesubsection}{1em}{}
\titlespacing*{\section}{\ParagraphIndent}{0pt}{\baselineskip}
\titlespacing*{\subsection}{\ParagraphIndent}{\baselineskip}{\baselineskip}
\usepackage{enumitem}
\setlist[itemize]{label=—,leftmargin=\ParagraphIndent,itemindent=0pt,nosep}
\setlist[enumerate]{label=\arabic*),leftmargin=\ParagraphIndent,nosep}
\usepackage[hidelinks,unicode]{hyperref}
\pagestyle{plain}
\emergencystretch=2em
\clubpenalty=10000
\widowpenalty=10000

% --------------------- ЗАПОЛНИТЕ ЭТИ ПОЛЯ ---------------------
% Наименование организации воспроизведено из предоставленного образца.
\newcommand{\University}{«САНКТ-ПЕТЕРБУРГСКИЙ НАЦИОНАЛЬНЫЙ
ИССЛЕДОВАТЕЛЬСКИЙ УНИВЕРСИТЕТ ИНФОРМАЦИОННЫХ
ТЕХНОЛОГИЙ, МЕХАНИКИ И ОПТИКИ»}
\newcommand{\Faculty}{Факультет безопасности информационных технологий}
\newcommand{\Discipline}{Название дисциплины}
\newcommand{\LabNumber}{1}
\newcommand{\LabTitle}{Название лабораторной работы}
\newcommand{\Group}{Номер группы}
\newcommand{\Student}{Фамилия Имя Отчество}
\newcommand{\Teacher}{Фамилия Имя Отчество}
\newcommand{\City}{Санкт-Петербург}
\newcommand{\ReportYear}{2026}

% Нумерация титульного листа учитывается, но номер не печатается.
\newcommand{\LabSection}[1]{%
  \ifNewSectionPage\clearpage\fi\section{#1}}
\newcommand{\StructuralSection}[1]{%
  \clearpage
  \phantomsection\addcontentsline{toc}{section}{#1}%
  {\centering\bfseries #1\par}%
  \nobreak\vspace{\baselineskip}}

\begin{document}
\begin{titlepage}
\thispagestyle{empty}
\setstretch{1.0}
\setlength{\parindent}{0pt}
\begin{center}
ФЕДЕРАЛЬНОЕ ГОСУДАРСТВЕННОЕ АВТОНОМНОЕ
ОБРАЗОВАТЕЛЬНОЕ УЧРЕЖДЕНИЕ ВЫСШЕГО ОБРАЗОВАНИЯ\par
\University\par
\vspace{10mm}
\Faculty\par
\vfill
Дисциплина:\par
«\Discipline»\par
\vspace{12mm}
\textbf{ОТЧЕТ ПО ЛАБОРАТОРНОЙ РАБОТЕ №\,\LabNumber}\par
\LabTitle\par
\vfill
\end{center}
\begin{flushright}
\begin{minipage}{0.55\textwidth}
\raggedleft
Выполнил:\par
Студент гр. \Group\par
\Student\par
Подпись: \rule{35mm}{0.4pt}\par
\vspace{8mm}
Проверил:\par
\Teacher\par
Подпись: \rule{35mm}{0.4pt}\par
Отметка: \rule{35mm}{0.4pt}\par
\end{minipage}
\end{flushright}
\vfill
\begin{center}
\City\par
\ReportYear\,г.
\end{center}
\end{titlepage}
\setcounter{page}{2}

% Содержание включайте, если оно требуется преподавателем:
% \renewcommand{\contentsname}{СОДЕРЖАНИЕ}
% \tableofcontents

\LabSection{Цель работы}
Сформулируйте цель работы: какой метод нужно изучить, какую зависимость
исследовать или какой практический навык получить.

Укажите исходные данные, номер варианта и поставленные задачи.

\LabSection{Ход работы}
\subsection{Оборудование и исходные данные}
Перечислите оборудование, программное обеспечение с версиями и параметры
эксперимента, необходимые для повторения работы.

\subsection{Выполнение и результаты}
Опишите выполненные действия и полученные результаты. Пример оформления
таблицы приведён в таблице~\ref{tab:results}. Замените демонстрационные
значения собственными данными.

\begin{table}[htbp]
\caption{Результаты измерений}\label{tab:results}
\centering
\begin{tabularx}{\textwidth}{|>{\centering\arraybackslash}p{22mm}|*{2}{>{\centering\arraybackslash}X|}}
\hline
Номер опыта & Напряжение $U$, В & Сила тока $I$, А \\ \hline
1 & 2,0 & 0,10 \\ \hline
2 & 4,0 & 0,20 \\ \hline
\end{tabularx}
\end{table}
\par\vspace{\baselineskip}

Расчёт сопротивления выполняется по формуле~\eqref{eq:resistance}:
\begin{equation}
R=\frac{U}{I},\label{eq:resistance}
\end{equation}
\noindent где $R$ — сопротивление, Ом; $U$ — напряжение, В;
$I$ — сила тока, А.

Пример размещения иллюстрации показан на рисунке~\ref{fig:example}.
\begin{figure}[htbp]
\centering
% Замените рамку командой: \includegraphics[width=0.8\textwidth]{image.png}
\fbox{\parbox[c][25mm][c]{0.75\textwidth}{\centering Место для рисунка}}
\caption{Название иллюстрации}\label{fig:example}
\end{figure}

\StructuralSection{ЗАКЛЮЧЕНИЕ}
Сформулируйте выводы по результатам работы. Укажите, достигнута ли цель,
какие значения и зависимости получены, чем объясняются отклонения
и какие ограничения есть у использованного метода.

% Если использованы внешние материалы, добавьте список и ссылки [1] в тексте:
% \StructuralSection{СПИСОК ИСПОЛЬЗОВАННЫХ ИСТОЧНИКОВ}
% \begin{enumerate}[label=\arabic*.]
% \item Автор. Название. — Город: Издательство, год. — Число страниц.
% \item Название ресурса. — URL: \url{https://example.org}
% (дата обращения: ДД.ММ.ГГГГ).
% \end{enumerate}

% При необходимости добавьте приложение с новой страницы:
% \StructuralSection{ПРИЛОЖЕНИЕ А}
% \begin{center}Дополнительные материалы\end{center}
% Здесь размещаются исходные данные, листинги или большие таблицы.
\end{document}
